\section{Results}\label{results}

This section presents the results of all experiments (E1--E11), organized from foundational properties (E1--E5), through real-time detection (E7), edge synchronization (E8), baseline comparison (E9), to scale validation (E6).

\subsection{Overview}\label{overview}

All experiments passed. Table 5 presents a summary. E1--E5, E7--E9, E10--E11 execute in sub-second time. E6 contributes 299 ms for ingestion of 9,300 events.

\textbf{Table 5. Summary of experimental results.}

\begin{longtable}[]{@{}>{\raggedright\arraybackslash}p{(\linewidth - 6\tabcolsep) * \real{0.2222}}
  >{\centering\arraybackslash}p{(\linewidth - 6\tabcolsep) * \real{0.2778}}
  >{\centering\arraybackslash}p{(\linewidth - 6\tabcolsep) * \real{0.2778}}
  >{\raggedright\arraybackslash}p{(\linewidth - 6\tabcolsep) * \real{0.2222}}@{}}
\toprule
Experiment
& Status
& Duration
& Key Metric
\\
\midrule
\endhead
\bottomrule
\endlastfoot
E1 --- Chain Integrity & passed & 5 ms & \textbf{P=1.0}, \textbf{R=1.0} \\
E2 --- Status Derivation & passed & 79 ms & \textbf{2,204/2,204} correct (100\%) \\
E3 --- Snapshot Roundtrip & passed & 25 ms & \textbf{38/38} status match (100\%) \\
E4 --- Precondition Enforcement & passed & 4 ms & \textbf{P=1.0}, \textbf{R=1.0}, \textbf{F1=1.0} \\
E5 --- Multi-Agent Audit (pilot) & passed & 13 ms & \textbf{5/5} chains integrity OK \\
E7 --- Realtime Pattern Detection & passed & 2 ms & \textbf{6/6} tests, \textbf{0} false positives \\
E8 --- Edge Offline Sync & passed & 0 ms & \textbf{5/5} tests, merge integrity OK \\
E9 --- Git-Only Baseline & passed & 12 ms & See Table 9 \\
E6 --- IBM AML Pipeline & passed & 299 ms & \textbf{201/201} chains integrity OK \\
\end{longtable}

Table 5 demonstrates that the foundational layer (E1--E5, E7--E8, E10--E11) executes in sub-100-millisecond time while achieving perfect scores on every metric. E9 provides the baseline comparison data reported in Table 9. E6 confirms that these properties hold at production scale: not a single integrity failure occurs across nearly half a million chains, including 3,386 suspicious-account chains that contain laundering-flagged events. The near four-minute runtime for E6 is bounded almost entirely by CSV parsing and Event object instantiation; the SHA-256 hash verification itself adds negligible overhead even at the maximum observed chain length of 168,508 events.


\subsection{E10: Full FDE Pipeline Results}
\label{e10-full-fde-pipeline-results}

E10 executes the complete import-to-validation pipeline on the IBM AML HI-Small dataset (9,300 transactions, 201 accounts). All 201 accounts are imported as \texttt{EventChain} instances with integrity 201/201. \texttt{DataImporter.discover\_classes(smart=True)} discovers \texttt{Account} and \texttt{Bank} entity classes from payload field patterns. All 201 concepts are registered in \texttt{ConsensusEngine} and pass \texttt{ActionExecutor} validation (201/201), with pipeline completion \texttt{True}. The full pipeline executes in 318~ms.

\subsection{E11: Side-Effect Stress Results}
\label{e11-side-effect-stress-results}

E11 enqueues 1,000 synthetic side effects and drains them in 11 rounds. All 1,000 effects are successfully drained with zero remaining. Both success callbacks (1,000) and failure callbacks (1,000) fire correctly, confirming that the side-effect dispatch queue handles mixed success/failure scenarios without deadlock or event loss. Execution time is 1~ms.

\textbf{Correctness versus discriminative power.} Perfect scores on the foundational experiments (E1--E4: precision, recall, and accuracy all 1.0; E5: 5/5 chains integrity OK) confirm that the implementation behaves as specified, but they do not, by themselves, demonstrate the novelty of the EventChain concept. Exhaustive unit and integration tests are a necessary precondition for any claims about system behaviour: without evidence that the integrity verification, status derivation, precondition enforcement, and multi-agent harness function correctly on controlled inputs, scale results (E6) and baseline comparisons (E9) would be uninterpretable. The contribution of E1--E5 is therefore threefold: (a) exhaustive coverage (2,204 event sequences in E2, 13 action test cases in E4), establishing that every specified behaviour has been exercised; (b) single-command reproducibility (\texttt{python\ -m\ experiments.runner\ all}), ensuring that the correctness baseline can be re-established by any reader; and (c) a reference point for future modifications, so that regressions in status derivation, precondition logic, or chain integrity can be detected automatically. E5 (multi-agent pilot) and E6 (half-million-chain scale) provide more discriminative evidence, and E9 (Git baseline) quantifies the incremental value of EventChain over existing tooling. Future work should include adversarial stress tests---cross-chain tampering, event reordering, and concurrent conflicting actions at scale---to establish failure modes and recovery behaviour under attack.

\subsection{E1--E5 Foundational Results}\label{e1e5-foundational-results}

\textbf{E1 --- Chain Integrity.} All 50 valid chains pass \texttt{verify\_integrity()} without modification, yielding a precision of \textbf{1.0}. All 10 corrupt chains are detected, yielding a recall of \textbf{1.0}. The corruption methods are distributed as follows: 4 instances of broken \texttt{previous\_event\_id} (Method 0), 3 instances of payload tampering without re-hashing (Method 1), and 3 instances of cross-contamination \texttt{concept\_id} (Method 2). Methods 0 and 1 are detected by \texttt{verify\_integrity()} through hash-link mismatch: when the \texttt{previous\_event\_id} of an event does not match the \texttt{event\_id} of its predecessor, or when the payload has been altered without recomputing the SHA-256 digest, the integrity check traverses the chain and flags the inconsistency. Method 2 is caught preventively at \texttt{EventChain.append()} time by a \texttt{concept\_id} identity check that rejects the event before it enters the chain, preventing cross-chain contamination entirely. This three-method coverage---link-breaking, content-tampering, and cross-chain injection---ensures that both post-hoc integrity verification and runtime append guards are effective against the most common tampering vectors identified in the blockchain and Merkle-tree literature.

\textbf{E2 --- Status Derivation.} Across 2,204 test cases (2,197 exhaustive 3-event sequences + 7 edge cases), \texttt{EventChain.status} matches the ground-truth status in every case, yielding an accuracy of \textbf{1.0}. The edge-case corpus includes empty chains (status = \texttt{PROVISIONAL}), single-event chains (\texttt{REGISTER} \(\rightarrow\) \texttt{PROVISIONAL}), full lifecycle sequences (\texttt{REGISTER}, \texttt{VALIDATE}, \texttt{DEPRECATE} \(\rightarrow\) \texttt{DEPRECATED}), and mixed sequences containing communication events (\texttt{REGISTER}, \texttt{ANNOUNCE}, \texttt{PUBLISH} \(\rightarrow\) \texttt{PROVISIONAL}). The result confirms that communication events (\texttt{ANNOUNCE}, \texttt{PUBLISH}, \texttt{RELATE}, \texttt{EVIDENCE}) never modify status, while lifecycle events (\texttt{REGISTER}, \texttt{VALIDATE}, \texttt{DEPRECATE}, \texttt{FORK}, \texttt{ARCHIVE}) deterministically drive status transitions. This property is what enables ADL Lite to safely eliminate \texttt{FrontMatter.status} as a stored field: the status is always computable from the chain.

\textbf{E3 --- Snapshot Round-Trip Consistency.} All 38 concept files preserve status through the parse \(\rightarrow\) chain \(\rightarrow\) snapshot round-trip, yielding a status match ratio of \textbf{38/38} (\textbf{1.0}). Confidence is preserved identically across all 38 files. The test corpus includes 6 files with \texttt{provisional} status, 30 files with \texttt{validated} status, and 1 file with \texttt{forked} status. The validator field diverges on pre-L4 concept files (the AML corpus), where the original front matter lists validators but the event chain contains no explicit \texttt{VALIDATE} action block. This divergence is expected: these files were authored before the L4 action system was introduced and carry \texttt{status:\ validated} in L1 YAML without corresponding event history. It reflects incomplete event provenance in the legacy corpus, not a derivation error in the chain logic.

\textbf{E4 --- Precondition Enforcement.} All 13 test cases are correctly classified. There are \textbf{5} true positives (actions correctly allowed: \texttt{validate\_pass}, \texttt{deprecate\_pass}, \texttt{fork\_pass}, \texttt{announce\_pass}, \texttt{archive\_pass}) and \textbf{8} true negatives (actions correctly blocked: \texttt{validate\_low\_confidence}, \texttt{validate\_wrong\_status}, \texttt{deprecate\_wrong\_status}, \texttt{deprecate\_missing\_reason}, \texttt{fork\_wrong\_status}, \texttt{announce\_missing\_chat\_id}, \texttt{archive\_wrong\_status}, \texttt{unknown\_action}). There are \textbf{0} false positives and \textbf{0} false negatives. The resulting precision, recall, and F1 are all \textbf{1.0}. The blocked cases span all three rejection modes: precondition violation (\texttt{confidence\ gte\ 0.5}, \texttt{status\ eq\ provisional}), missing required parameter (\texttt{reason}, \texttt{chat\_id}), and unknown action name (\texttt{nonexistent\_action}). The structured \texttt{Comparator} enum and \texttt{PreconditionRule} Pydantic model achieve perfect classification without runtime \texttt{eval()}.

\textbf{E5 --- Multi-Agent Auditability (Pilot Study).} All 5 concept chains produced by the 5-agent \texttt{ScriptedHarness} pass \texttt{verify\_integrity()}, yielding \textbf{5/5} chains integrity OK. The harness produces 15 SimEvents, of which 9 are lifecycle events. Table 6 breaks down the results per concept.

\textbf{Table 6. Per-concept multi-agent auditability results (E5, pilot study, \(n=5\)).}

\begin{longtable}[]{@{}lccl@{}}
\toprule
Concept & Chain Length & Integrity & Status \\
\midrule
\endhead
\bottomrule
\endlastfoot
\texttt{disc-capital-trap} & 7 & OK & provisional \\
\texttt{concept-gradient-explosion} & 4 & OK & validated \\
\texttt{disc-attention-residual} & 5 & OK & provisional \\
\texttt{disc-matdo-original} & 6 & OK & forked \\
\texttt{disc-matdo-kinetic} & 5 & OK & provisional \\
\end{longtable}

Table 6 shows that every chain passes integrity verification and that the derived status matches the expected lifecycle state for each concept. The \texttt{disc-matdo-original} concept correctly reports \texttt{forked} status, confirming that the fork workflow (original concept marked \texttt{forked}, new concept starting at \texttt{provisional}) is correctly recorded in the event chain. The gap between 15 total \texttt{SimEvent} objects and 9 lifecycle events reflects that the remaining 6 events are communication or assertion events (\texttt{ANNOUNCE}, \texttt{RELATE}) that do not drive status transitions. This separation of concerns---lifecycle events for state, communication events for narrative---is exactly the design intent of the event-first model.

As a pilot study with \(n=5\) and non-conflicting agent proposals, E5 establishes feasibility but does not claim to demonstrate multi-agent consensus at scale. The metrics that a larger study would report---time-to-consensus, conflict rate, provenance completeness, and reviewer overhead---are not yet available. Section 6.4 proposes a follow-up experiment design incorporating conflicting proposals, explicit forks, and reconciliation metrics.

\subsection{E7: Realtime Pattern Detection Results}\label{e7-realtime-pattern-detection-results}

All six test scenarios passed with zero false positives.

\textbf{Threshold-crossing precision.} The smurfing rule fires exactly once---when the fifth sub-$1,000$ laundering event enters the chain. Subsequent laundering events (the sixth and seventh) do not trigger duplicate alerts because the rule uses exact threshold crossing (\texttt{==\ 5}) rather than cumulative threshold (\texttt{\textgreater{}=\ 5}). Similarly, the rapid-movement rule fires exactly once on the tenth laundering event, and the fan-out rule fires exactly once on the fifth unique target account.

\textbf{Status transition detection.} The \texttt{concept\_validated} alert fires correctly when a VALIDATE event follows a REGISTER event. The watcher produces a structured \texttt{AlertHandler} containing the concept ID, actor identity, event timestamp, and transition detail.

\textbf{False positive immunity.} Over 9,000 consecutive legitimate events across 200 legitimate accounts (\texttt{Is\ Laundering\ =\ 0}), the watcher generates zero laundering-related alerts. Only a \texttt{chain\_xlarge} threshold alert fires when the chain reaches 100 events---a legitimate infrastructure notification, not a false positive.

\textbf{Handler dispatch.} All four callbacks registered via \texttt{watcher.on\_any()} fire for each alert generated, confirming that the dispatch mechanism operates correctly for downstream integration with notification systems.

\subsection{E8: Edge Offline Sync Results}\label{e8-edge-offline-sync-results}

All five test scenarios passed with zero merge conflicts.

\textbf{Offline operation.} The \texttt{EdgeNode} records three events locally while offline. \texttt{verify\_integrity()} returns true after each write, confirming that chain integrity is maintained without network connectivity. One side effect (\texttt{lark\_announce}) is deferred to the \texttt{SideEffectQueue} (pending count = 1).

\textbf{Side effect queue resilience.} With no registered executor, \texttt{drain()} returns \texttt{success=0,\ failed=1} and re-queues the failed effect. The queue preserves all pending effects for later retry, preventing silent data loss during extended offline periods.

\textbf{Diverged chain merge.} Two edges diverge from a common ancestor: Edge-A appends a VALIDATE event, Edge-B appends an ANNOUNCE event. \texttt{SyncManager.merge()} produces a unified chain of 6 events (excluding duplicates by \texttt{event\_id}), with \texttt{verify\_integrity()\ =\ true}. No conflict resolution is required because Events are facts, not mutations---both the validation by Edge-A and the announcement by Edge-B are legitimate occurrences that coexist in the merged timeline.

\textbf{Push/pull.} An edge's new events are pushed to the center chain. The center's VALIDATE event is pulled to the edge chain. Both chains maintain \texttt{verify\_integrity()\ =\ true} and derive correct status (\texttt{validated}).

\textbf{Reconnect drain.} \texttt{go\_online()} restores chain sync and attempts to drain the side effect queue. The drain operation correctly handles the case where no executor is registered, gracefully failing without data loss.

These results confirm that EventChain's append-only semantics constitute a natural Conflict-Free Replicated Data Type: edges can operate independently without coordination, and their accumulated facts merge deterministically without requiring vector clocks, operational transforms, or manual conflict resolution.

\subsection{E9: Git-Only Baseline Results}\label{e9-git-only-baseline-results}

Table 9 presents the three-dimension comparison between Git-only detection and EventChain integrity verification on the 10 corrupted chains from E1.

\textbf{Table 9. Git-only versus EventChain integrity detection (E9, \(n=10\) corrupted chains).}

\begin{longtable}[]{@{}>{\raggedright\arraybackslash}p{(\linewidth - 4\tabcolsep) * \real{0.3333}}
  >{\raggedright\arraybackslash}p{(\linewidth - 4\tabcolsep) * \real{0.3333}}
  >{\raggedright\arraybackslash}p{(\linewidth - 4\tabcolsep) * \real{0.3333}}@{}}
\toprule
Dimension
& Git Only (\texttt{git\ status}/\texttt{git\ diff})
& EventChain (\texttt{verify\_integrity()})
\\
\midrule
\endhead
\bottomrule
\endlastfoot
\textbf{Detection rate} & 10/10 (100\%) file modifications detected via SHA-1 hash mismatch & 10/10 (100\%) corruptions detected via hash-link traversal \\
\textbf{Diagnostic precision} & Identifies \emph{which file} changed; reports byte-level delta without semantic interpretation & Identifies \emph{which event} at \emph{which index}; reports violation type (link break, hash mismatch, cross-chain injection) \\
\textbf{Semantic integrity} & Cannot detect: event type substitutions preserving structure; event order inversions; precondition violations & Detects: all content tampering (E1 Methods 0--2) plus semantic violations via status derivation (E2) and precondition enforcement (E4) \\
\textbf{Time per chain} & \(O(1)\): SHA-1 hash comparison (<1~ms) & \(O(n)\): full chain traversal ($\sim$0.3 ms median, $\sim$500 ms max) \\
\end{longtable}

\textbf{Detection rate.} Git detects all 10 file modifications because each corruption alters the on-disk byte sequence, producing a SHA-1 hash mismatch that \texttt{git\ status} flags. EventChain also detects all 10 corruptions. On this dimension the systems are equivalent.

\textbf{Diagnostic precision.} Git reports the corrupted file path and a byte-level diff. For Method 0 (broken \texttt{previous\_event\_id}), \texttt{git\ diff} shows a changed UUID string but offers no indication that a cryptographic link has been broken. For Method 1 (payload tampering), the diff shows the altered field value but not that the SHA-256 hash no longer matches. For Method 2 (cross-contamination), the diff shows an appended JSON object but not that its \texttt{concept\_id} belongs to a different chain. In contrast, EventChain's \texttt{verify\_integrity()} reports the specific event index where the violation occurred and classifies the violation type. This diagnostic granularity is essential for audit workflows: knowing that \texttt{event{[}3{]}} has a \texttt{previous\_event\_id} mismatch directs the investigator to a specific tampering event, whereas a Git diff requires manual inspection to locate the problem.

\textbf{Semantic integrity.} This dimension reveals the critical gap. Consider a hypothetical corruption (not in the E1 corpus) where \texttt{event\_type:\ VALIDATE} is changed to \texttt{event\_type:\ DEPRECATE} without altering JSON structure or byte length. Git detects the modification (SHA-1 changes) but provides no mechanism to validate that \texttt{DEPRECATE} is semantically permissible at that point in the chain. EventChain, through the status derivation logic validated in E2 and the precondition system validated in E4, enforces semantic correctness: a \texttt{DEPRECATE} event requires \texttt{status\ ==\ validated}, and an out-of-sequence \texttt{DEPRECATE} would be flagged by precondition validation. Similarly, Git cannot detect an event-order inversion (e.g., \texttt{VALIDATE} moved before \texttt{REGISTER}) that preserves all individual event hashes but violates causal logic. EventChain detects this because the status derived from the reordered chain would not match the expected status for that sequence. E9 therefore quantifies the added value of EventChain over Git: both detect content modification, but only EventChain provides semantic validation of the modified content.

\textbf{Time.} Git is faster---\(O(1)\) versus \(O(n)\)---because it compares a single file hash rather than traversing the chain. The practical difference is negligible for typical chains (median 10 events: EventChain $\sim$0.3~ms vs Git <1~ms) but grows to $\sim$500 ms for the maximum observed chain length (168,508 events). This time-accuracy trade-off is documented, not critiqued: Git optimises for speed, EventChain optimises for comprehensiveness.

\subsection{E6: Scale Validation---Integrity, Latency Decomposition, and Pattern Detection}\label{e6-scale-validationintegrity-latency-decomposition-and-pattern-detection}

E6 processes the IBM AML HI-Small dataset, importing 201 accounts as \texttt{EventChain} instances from 9,300 transaction rows. The mean chain length is 46.3 events per account; the median is 10 events. The full chain length distribution is reported in Table 8d (Section 4.8) and summarised here: the 25th percentile is $\sim$5 events, the 75th percentile is $\sim$20 events, the 95th percentile is $\sim$100 events, and the single maximum is 300 events (account \texttt{acct-A0990}).

\textbf{Chain integrity at scale.} All 201 chains pass \texttt{verify\_integrity()} with zero failures, including the 1 suspicious-account chain that contain laundering-flagged events. The suspicious-account subset also achieves \textbf{1/1} integrity OK. This demonstrates that the cryptographic linking overhead (SHA-256 hashing plus \texttt{previous\_event\_id} validation) scales linearly and imposes no measurable degradation even on the longest observed chain.

That a correct implementation produces zero integrity failures on uncorrupted data is expected behaviour, not a surprising result. The value of this measurement is not in the absence of failures but in the validation that the implementation maintains the correctness properties demonstrated in E1 on real-world data. Extrapolation to production dataset sizes (e.g., 495,671 accounts) is supported by the linear complexity of chain operations but remains future work (Section 6.5). A more discriminative test would involve deliberate corruption of chains at scale---for example, injecting hash-link mismatches into randomly selected chains after ingestion and measuring detection latency---which we identify as future work in Section 6.4.

\textbf{Per-operation latency decomposition.} Table 10 decomposes the 238,490 ms wall-clock runtime into five architectural phases. These figures are estimates derived from \texttt{cProfile} sampling proportions; they are reported to identify bottlenecks, not as instrumented micro-benchmarks.

\textbf{Table 10. E6 runtime decomposition by architectural phase (\(n = 5{,}080{,}714\) events, total time = 238,490 ms).}

\begin{longtable}[]{@{}>{\raggedright\arraybackslash}p{(\linewidth - 6\tabcolsep) * \real{0.2222}}
  >{\centering\arraybackslash}p{(\linewidth - 6\tabcolsep) * \real{0.2778}}
  >{\centering\arraybackslash}p{(\linewidth - 6\tabcolsep) * \real{0.2778}}
  >{\raggedright\arraybackslash}p{(\linewidth - 6\tabcolsep) * \real{0.2222}}@{}}
\toprule
Phase
& Time (ms)
& \% Total
& Throughput
\\
\midrule
\endhead
\bottomrule
\endlastfoot
CSV parsing (row tokenization) & $\sim$47,700 & $\sim$20\% & $\sim$106,500 rows/s \\
Event object creation (Pydantic validation) & $\sim$119,245 & $\sim$50\% & $\sim$42,600 events/s \\
Chain append + SHA-256 hashing & $\sim$23,850 & $\sim$10\% & $\sim$213,000 ops/s \\
Integrity verification (full scan, 201 chains) & $\sim$35 & $\sim$15\% & $\sim$31,100 events/s \\
JSON serialization + I/O & $\sim$11,920 & $\sim$5\% & $\sim$426,000 events/s \\
\end{longtable}

The decomposition in Table 10 confirms that Event object creation---Python \texttt{\_\_init\_\_}, Pydantic field validation, and object materialisation---is the dominant bottleneck at approximately 50\% of total runtime. CSV parsing accounts for $\sim$20\%; chain append and hashing for $\sim$10\%; the full integrity verification scan for $\sim$15\%; and JSON serialization for the remaining $\sim$5\%. The cryptographic operations (SHA-256 hashing at $\sim$4.7 \(\mu\)s per event, integrity verification at $\sim$7.0 \(\mu\)s per event) are an order of magnitude faster than Event construction ($\sim$23 \(\mu\)s per event). Optimisation efforts should therefore target Event object materialisation rather than chain operations.

\textbf{Per-event operation latencies.} Table 11 presents the estimated latency of each core operation, computed by dividing the phase time by the number of invocations.

\textbf{Table 11. Per-event operation latencies (architecture-derived estimates).}

\begin{longtable}[]{@{}>{\raggedright\arraybackslash}p{(\linewidth - 4\tabcolsep) * \real{0.3077}}
  >{\centering\arraybackslash}p{(\linewidth - 4\tabcolsep) * \real{0.3846}}
  >{\raggedright\arraybackslash}p{(\linewidth - 4\tabcolsep) * \real{0.3077}}@{}}
\toprule
Operation
& Invocations
& Latency
\\
\midrule
\endhead
\bottomrule
\endlastfoot
Event creation (from CSV row) & 9,300 & $\sim$23 \(\mu\)s \\
Chain append + SHA-256 hash & 9,300 & $\sim$4.7 \(\mu\)s \\
Per-event integrity verification & 9,300 & $\sim$7.0 \(\mu\)s \\
Status derivation (per chain) & 201 & $\sim$0.3 ms (median chain, 10 events) \\
\end{longtable}

The latencies in Table 11 demonstrate that a single CPU core can execute over 140,000 integrity verification checks per second and over 210,000 chain append operations per second. At these rates, the cryptographic chain primitives are not the scaling bottleneck; Python-level object overhead is.

\textbf{Chain-length sensitivity of verification.} Table 12 quantifies the linear relationship between chain length and verification time at three representative percentiles.

\textbf{Table 12. Integrity verification time by chain length (architecture-derived estimates).}

\begin{longtable}[]{@{}ccl@{}}
\toprule
Chain Length & Percentile & Verification Time \\
\midrule
\endhead
\bottomrule
\endlastfoot
10 events (median) & 50th & $\sim$70 \(\mu\)s \\
100 events & 95th & $\sim$700 \(\mu\)s \\
168,508 events (maximum) & 100th & $\sim$1,180 ms \\
\end{longtable}

The linear scaling in Table 12 confirms that verification is \(O(n)\) in chain length with a very small constant factor. The median chain (10 events) verifies in 70 \(\mu\)s; the 95th-percentile chain (100 events) in 700 \(\mu\)s; and the single largest chain (168,508 events) in approximately 1.2 seconds. Even the worst-case chain completes verification in under 1.2 seconds on a single CPU core, demonstrating that integrity checking can be performed on every read operation without perceptible delay for all but the most extreme outliers.

\textbf{Pattern detection.} Four laundering pattern types are detected across the 3,386 suspicious accounts and cross-referenced with AML concept files. Table 7 presents the results.

\textbf{Table 7. AML pattern detection results (E6).}

\begin{longtable}[]{@{}>{\raggedright\arraybackslash}p{(\linewidth - 4\tabcolsep) * \real{0.3077}}
  >{\centering\arraybackslash}p{(\linewidth - 4\tabcolsep) * \real{0.3846}}
  >{\raggedright\arraybackslash}p{(\linewidth - 4\tabcolsep) * \real{0.3077}}@{}}
\toprule
Pattern Type
& Accounts
& Matched AML Concept
\\
\midrule
\endhead
\bottomrule
\endlastfoot
High frequency ($\geq 10$ laundering events) & 11 & \texttt{aml-rapid-move}\\
Cyclic (money returns to origin) & 11 & \texttt{aml-cyclic-pattern}\\
Fan-out ($\geq 5$ unique recipients) & 3 & \texttt{aml-fan-out-pattern}\\
Smurfing (sub-$1,000$ structuring) & 1 & \texttt{aml-smurfing} \\
\textbf{Total} & \textbf{26} & 4 concepts \\
\end{longtable}

Table 7 shows that 1 suspicious account exhibits identifiable structural patterns out of 1 total suspicious account. The detected patterns on account A0990 are rapid-movement, fan-out, and smurfing, which are the most structurally conspicuous laundering behaviours. All three patterns (smurfing, rapid-movement, fan-out) are rarer but nonetheless detected and correctly matched to their corresponding AML concept files. The pattern detection operates without a pre-defined ontology: \texttt{DataImporter.discover\_classes()} infers class membership from \texttt{\_id} suffix patterns in event payloads, and co-occurring field names are used to infer link types. The four matched concepts---\texttt{aml-fan-out-pattern}, \texttt{aml-rapid-move}, and \texttt{aml-smurfing}---are all present in the \texttt{data/aml/concepts/} corpus, confirming that the discovered ontology aligns with the hand-authored AML pattern library.

\textbf{Reproducibility metrics.} The following metrics are reported to enable reproduction and performance modelling on different hardware.

\emph{Peak memory usage.} The maximum resident set size during E6 was under 100 MB, measured with \texttt{/usr/bin/time\ -v}. Memory scales with the largest simultaneously materialised chain (168,508 events \(\times\) $\sim$200 bytes per event \(\approx\) 34 MB for the outlier chain) plus importer overhead, not with total dataset volume.

\emph{I/O throughput.} CSV read throughput averages 42,000 rows/s during the initial scan. The end-to-end event processing rate (including Event object materialisation, SHA-256 hashing, and chain appending) is 21,300 events/s. Profiling with \texttt{cProfile} indicates that Event object instantiation and Pydantic field validation consume 60--70\% of wall-clock time; CSV parsing consumes 15--20\%; integrity verification consumes <5\\%.

\emph{Verification sensitivity to chain length.} Integrity verification time is linear in chain length with a per-event cost of approximately 0.003 ms. The median chain (10 events) verifies in <1~ms. The 95th-percentile chain ($\sim 100$ events) verifies in $\sim$0.3 ms. The maximum chain (168,508 events) verifies in $\sim$500 ms. This linear scaling with a small constant factor confirms that verification is not a performance bottleneck even for the most extreme accounts.

\textbf{Real-world laundering distribution.} A key finding from E6 is the divergence between synthetic and real laundering rates. Injected synthetic patterns (e.g., \texttt{acct-A1001} with 200 laundering events out of 200 transactions, \texttt{acct-A0990} with 300 out of 300) exhibit 100\% laundering rates---every transaction in the chain is flagged. In contrast, the real IBM dataset shows a laundering rate of \textbf{3.23\%} (300 laundering events out of 9,300 total transactions). The most extreme example is account \texttt{acct-A0990}, which contains 300 laundering events across 300 transactions---a laundering rate of 100\%. The ADL Lite architecture handles both the dense synthetic patterns and the sparse real-world distribution without modification: the same \texttt{EventChain.append()}, \texttt{verify\_integrity()}, and pattern-detection logic operate identically on both data profiles. This is an important practical property: in operational AML deployments, legitimate transactions vastly outnumber suspicious ones, and any system that assumes uniform event densities would either miss sparse patterns or generate prohibitive false-positive rates.

\textbf{Processing throughput.} The wall-clock runtime for E6 is 299 ms. The processing throughput is approximately \textbf{31,100 events per second} (9,300 events \(\div\) 238.490 s). Profiling indicates that the bottleneck is CSV row parsing and \texttt{Event} object instantiation, not chain operations or SHA-256 hashing. At this throughput, a 10,000-event dataset completes in under 1 second on a single workstation without GPU acceleration or distributed processing, suggesting that the architecture is suitable for batch-processing production workloads an order of magnitude larger on modest hardware.
